\begin{papercard}{schultchen2021}
\cf{Stages} \stagetag{A} \stagetag{D} \stagetag{I}. A German-led systematic store review of mindfulness apps in the European stores, using a web crawler and dual MARS-G rating with ICC. The closest German-market precedent for store search mechanics.

\cf{Aim} Systematically review and rate mindfulness apps in the European stores, with a broader search than earlier single-term, single-store reviews, plus privacy policy features and the relation between user stars and expert quality \pc{schultchen2021}{2}{the aim in this study was to conduct a system}.

\cf{Search} An automated search engine (web crawler) queried the European Apple App Store and Google Play on 24 October 2018 with four keywords: ``mindfulness'', ``meditation'', ``yoga'', ``body scan'' \pc{schultchen2021}{2}{Using an automated search engine (webcrawler), a search in}. Results went into a central database, and duplicates were removed automatically \pc{schultchen2021}{2}{All apps were listed in a central database, and duplicates}. The crawler was built in the Mobile Health App Database (MHAD) project at Ulm \pc{schultchen2021}{6}{port in the development of the search engine and their support in the}.

\cf{Eligibility} (1) Sufficient content conceptualised for mindfulness, excluding fitness/yoga apps and apps offering only timers, reminders, music or quotes; (2) German or English, in line with reviewer language skills; (3) officially available in the European store. Second step: download and check, including full functionality (no device problems, not in testing) \pc{schultchen2021}{2}{inclusion criteria: (1) sufficient content and conceptualized}. Flow: 605 apps (434 iOS, 171 Android) screened, 343 excluded, 262 downloaded, 70 excluded (14 not mindfulness, 12 not working, 18 not available, 6 duplicates, 20 language or other), giving 192 included (123 English, 69 German) \pc{schultchen2021}{3}{The search resulted in 605 apps}.

\cf{Rating procedure} Two independent reviewers per app using the German MARS. All had a psychology degree and completed a free online MARS training video, and each app was tested for at least 15--20 minutes \pc{schultchen2021}{2}{Two independent reviewers evaluated each app using the}. Interrater reliability was computed as ICC, with $\geq$0.75 considered satisfactory. If an app's ratings fell below that, a third rater (DS or LS) was consulted, and final scores are the means of the two raters \pc{schultchen2021}{2}{The interrater reliability (IRR) between the two review}. Descriptive extraction: name, platform, content subcategories, price, goals, methods, data protection and privacy, store user rating, number of RCTs \pc{schultchen2021}{2}{To describe the app characteristics, different descriptive data}. Privacy and security features were rated according to the MARS classification on a 5-point scale \pc{schultchen2021}{2}{privacy and security features were reviewed}.

\cf{Results (method-relevant)} ICC = 0.87 (95\,\% CI 0.86--0.88) \pc{schultchen2021}{4}{The intra-class correlation, indicating agreement of both}. Only 23\,\% had a privacy policy: 7 with active confirmation and 12 with a ``passive'' policy. 8\,\% had password protection, and 24\,\% gave contact details or an imprint (\emph{Impressum}) \pc{schultchen2021}{4}{Data security and privacy features showed that 15 apps}. User star ratings correlated weakly with engagement, functionality and aesthetics but not with information quality or overall quality \pc{schultchen2021}{4}{No correlation could be found}.

\cf{Limitations stated} A fast-changing market, so a repeated search would differ. Only Apple and Google stores (about 99\,\% of the market). Search terms were limited. German/English only because of reviewer language. MARS only, no ENLIGHT or APA model \pc{schultchen2021}{6}{This study has several limitations}.

\cf{Use for this thesis} Precedent for a \emph{German/English language rule justified by reviewer competence}, a crawler-based store search with a recorded date, automatic de-duplication, a two-step screen (store page, then download), a PRISMA-style flow, and two raters with a pre-set ICC threshold and a third-rater rule. Note the domain rule: functional mindfulness content, excluding timer- or music-only apps. This mirrors the thesis' functional tracker definition. Its privacy coding (policy present, active vs.\ passive confirmation, imprint) is a simple German-market baseline.
\end{papercard}
